% !TeX program = xelatex
\documentclass[UTF8,zihao=-4]{ctexart}

% =========================================================
% 2026 APMCM 中文赛项 A题论文 LaTeX 模板
% 编译方式：XeLaTeX
% 注意：摘要页和正文不得出现队员姓名、学校、邮箱、学号等个人信息
% =========================================================

% -------------------- 页面与基础格式 --------------------
\usepackage[a4paper,left=2.5cm,right=2.5cm,top=2.5cm,bottom=2.5cm]{geometry}
\usepackage{setspace}
\usepackage{amsmath,amssymb,bm}
\usepackage{booktabs,longtable,array,multirow,tabularx}
\usepackage{graphicx}
\usepackage{float}
\usepackage{caption}
\usepackage{subcaption}
\usepackage{enumitem}
\usepackage{fancyhdr}
\usepackage{hyperref}
\usepackage{xcolor}
\usepackage{xurl}
\usepackage{listings}

% -------------------- 基本排版 --------------------
\setstretch{1.0}                 % 单倍行距
\setlength{\parindent}{2em}      % 首行缩进
\setlength{\parskip}{0pt}
\emergencystretch=2em            % 降低 overfull hbox 概率
\sloppy
\hypersetup{hidelinks}

% -------------------- 图表标题格式 --------------------
\captionsetup[figure]{name=图,labelsep=space,font=small}
\captionsetup[table]{name=表,labelsep=space,font=small}

% -------------------- 页眉页脚 --------------------
\pagestyle{fancy}
\fancyhf{}
\cfoot{\thepage}
\renewcommand{\headrulewidth}{0pt}

% -------------------- 表格列格式 --------------------
\newcolumntype{Y}{>{\centering\arraybackslash}X}
\newcolumntype{L}{>{\raggedright\arraybackslash}X}

% -------------------- 代码环境 --------------------
\lstset{
  basicstyle=\ttfamily\small,
  breaklines=true,
  frame=single,
  columns=fullflexible,
  numbers=left,
  numberstyle=\tiny,
  keywordstyle=\color{blue},
  commentstyle=\color{gray},
  stringstyle=\color{teal}
}

% -------------------- 自定义命令 --------------------
\newcommand{\teamid}{apmcm26203797} % 修改为正式队号
\newcommand{\problemid}{A}
\newcommand{\papertitle}{基于时滞识别与机理--数据融合的自来水厂水浊度预测及风险评估模型}

% =========================================================
\begin{document}

% =========================================================
% 摘要首页
% =========================================================

\begin{center}
\begin{tabularx}{0.95\textwidth}{|Y|Y|Y|Y|}
\hline
选题 & \problemid & 参赛编号 & \teamid \\
\hline
\end{tabularx}
\end{center}

\vspace{1.2em}

\begin{center}
{\heiti\zihao{3}\papertitle}
\end{center}

\vspace{1em}

\begin{center}
{\heiti\zihao{4} 摘要}
\end{center}

自来水厂出厂水水质直接关系到供水安全，其变化过程受到原水指标、操作变量等多种因素的共同影响。针对题目给出的多源时序监测数据，本文围绕影响因素识别、水质动态预测以及运行策略优化等问题展开研究，构建兼具解释性与预测能力的数学模型，为水厂精细化运行管理提供决策依据。

针对问题一，首先通过

针对问题二，

针对问题三，

针对问题四，


\vspace{1em}
\noindent\textbf{关键词：} 水浊度预测；时滞识别；机理--数据融合；随机森林；XGBoost；GRU；水质风险评价

\newpage

% =========================================================
% 正文
% =========================================================

\section{问题重述}

自来水厂生产过程包括取水、混凝、沉淀、过滤、消毒和出厂输送等多个环节。原水浊度、原水 pH、原水流量、矾投加量、滤后水浊度、清水池水位和出厂水流量等变量共同影响最终出厂水浊度。由于混凝反应、过滤过程和清水池调蓄作用均具有时间延迟，当前出厂水质并不完全由当前原水状态决定，而是受过去若干时刻水质和工艺参数共同影响。

本文需要基于给定的多变量时间序列数据，完成以下四项任务：
\begin{enumerate}[label=（\arabic*）]
  \item 筛选影响出厂水浊度 $NTU$ 的主要因素，解释各因素影响程度和作用方向，并预测 2026 年 2 月 1 日、2 月 10 日、2 月 20 日的水浊度；
  \item 建立滤后水浊度 $FILT.NTU$ 的动态时滞模型，描述原水指标和操作变量对滤后水浊度的滞后影响，并估计不同变量的时滞参数；
  \item 结合质量守恒思想和数据驱动方法，建立未来 1--12 小时出厂水浊度预测模型，并给出指定日期 7:00 至 19:00 的预测结果；
  \item 以 $NTU\leq 1$ 为硬性约束，结合超标幅度和异常持续时间，建立水质风险评价体系，对 2026 年近三个月水质进行风险分级。
  \item 在研究时间范围内，各输入变量对滤后水浊度（NTU）的时间滞后效应保持稳定。
\end{enumerate}

\section{问题分析}

本题本质上是一个具有多变量耦合、非线性响应和时间滞后的水处理过程建模问题。若直接采用普通静态回归模型，容易忽略矾投加、原水浊度变化和清水池调蓄对出厂水浊度的滞后影响；若仅采用深度学习模型，则解释性不足，难以说明不同变量的影响方向和作用机制。因此，本文采用“数据预处理--滞后特征构造--机器学习解释--动态预测--风险评价”的整体建模路线。

\begin{figure}[H]
  \centering
  \fbox{\parbox{0.88\textwidth}{\centering
  原始数据 $\rightarrow$ 时间重构与数据清洗 $\rightarrow$ 缺失值和异常值处理 $\rightarrow$ 滞后特征构造\\[0.5em]
  $\rightarrow$ 问题一：影响因素筛选与出厂水 NTU 预测\\[0.5em]
  $\rightarrow$ 问题二：$FILT.NTU$ 时滞动态模型\\[0.5em]
  $\rightarrow$ 问题三：机理--数据融合多步预测模型\\[0.5em]
  $\rightarrow$ 问题四：水质风险评价模型}}
  \caption{水浊度预测与风险评价总体建模流程}
  \label{fig:framework}
\end{figure}

\section{模型假设}

\begin{enumerate}[label=假设\arabic*：]
  \item 经过缺失值处理和异常值标记后的监测数据能够反映水厂主要运行规律。
  \item 原水水质和工艺操作变量对滤后水、出厂水浊度的影响存在有限时滞，且主要时滞范围不超过 24 小时。
  \item 清水池调蓄作用可由清水池水位和出厂水流量构造的近似水力停留时间特征表征。
  \item 短时传感器噪声不会改变水处理过程中的主要动态关系。
  \item 以生活饮用水浊度限值 $NTU\leq 1$ 作为判断出厂水是否超标的统一标准。
  \item 
\end{enumerate}

\section{符号说明}

\begin{table}[H]
\centering
\caption{主要符号说明}
\begin{tabular}{cl}
\toprule
符号 & 含义 \\
\midrule
$t$ & 时间序号 \\
$h$ & 预测步长 \\
$y_t$ & $t$ 时刻出厂水浊度 $NTU$ \\
$z_t$ & $t$ 时刻滤后水浊度 $FILT.NTU$ \\
$x_{i,t}$ & 第 $i$ 个输入变量在 $t$ 时刻的观测值 \\
$\tau_i$ & 第 $i$ 个变量的最优时滞 \\
$HRT_t$ & $t$ 时刻近似水力停留时间特征 \\
$E_t$ & 超标幅度 \\
$D_t$ & 连续异常持续时间 \\
$V_t$ & 浊度波动指标 \\
$R_t$ & 综合风险指数 \\
$RMSE$ & 均方根误差 \\
$MAE$ & 平均绝对误差 \\
$R^2$ & 决定系数 \\
\bottomrule
\end{tabular}
\end{table}

% =========================================================
\section{数据预处理与质量审计}

\subsection{数据合并与时间重构}
\subsubsection{数据合并}
原始附件由多个按月份或按日期工作表记录的 Excel 文件构成，不同文件在日期格式、字段完整性和缺失情况上存在差异。为保证后续时间序列分析与预测建模的可靠性，本文首先对所有原始表格进行统一读取与整合，最终得到包含 5460 条两小时观测记录的统一数据集。数据时间范围覆盖 2025 年 1 月 1 日 07:00 至 2026 年 4 月 1 日 05:00，共包含题目涉及的 21 个字段，合并后时间戳重复数为 0。
\subsubsection{时间重构}
由于部分原始附件同时存在 Excel 日期格式与“日/月/年”文本格式混用的问题，且个别日期可能被软件错误识别为“月/日”格式，直接使用原始日期列作为时间索引存在风险。因此，本文依据文件所属月份、每日固定观测顺序以及监测时间信息重建标准时间轴。

考虑到每日记录顺序为：

$$
07:00,09:00,\cdots,23:00,01:00,03:00,05:00
$$

本文将 01:00、03:00 和 05:00 三个时点归入下一自然日，从而构建严格递增的统一时间序列。

\subsection{变量标准化}
\subsubsection{字段统一}根据题目附录给出的变量定义，本文将不同来源文件中的表头统一映射为标准字段名称，包括原水流量、原水浊度、原水 pH、滤后水浊度、处理后水浊度、余氯、混凝剂投加量、出厂流量等变量。对于题目未明确给出物理单位的字段，本文不额外推断其量纲，仅保留原始数据形式。
\subsubsection{送水泵运行转台转换}
原始字段 T/W PUMP DUTY 使用泵组编号组合表示运行状态，例如“1\&3”、“2\&4”等。由于泵编号本身不具有数值大小意义，而运行泵数量能够更直接反映系统运行负荷，因此本文将该字段转换为运行泵数量特征。

具体地，通过提取记录中的泵号数字并统计其个数获得运行泵数量。例如：

\[1+2+5 \rightarrow 3\]

\[1,4 \rightarrow 2\]

该转换保留了运行强度信息，同时提高了变量的可解释性与模型适用性。
%=========================================================
\section{问题一：主要影响因素筛选与出厂水浊度预测模型}

\subsection{候选影响因素选取}

结合水处理工艺机理和数据字段，选取原水浊度、原水 pH、原水流量、滤后水浊度、矾投加量、清水池水位和出厂水流量等变量作为候选影响因素。

\begin{table}[H]
\centering
\caption{候选变量及物理含义}
\begin{tabularx}{0.95\textwidth}{l l L}
\toprule
变量 & 中文含义 & 建模作用 \\
\midrule
$R/W\ NTU$ & 原水浊度 & 主要输入变量，反映进入处理系统的颗粒污染负荷 \\
$R/W\ PH$ & 原水 pH & 影响混凝剂作用效果，可能具有非线性影响 \\
$R/W\ FLOW$ & 原水流量 & 影响处理负荷和水力停留状态 \\
$FILT.NTU$ & 滤后水浊度 & 反映过滤前后处理效果，是出厂水浊度的重要中间变量 \\
$ALUM$ & 矾投加量 & 工艺控制变量，对浊度变化具有滞后调节作用 \\
$C/W\ WELL\ LEVEL$ & 清水池水位 & 与清水池调蓄和水力停留过程相关 \\
$T/W\ FLOW$ & 出厂水流量 & 与出厂输送负荷和停留时间相关 \\
\bottomrule
\end{tabularx}
\end{table}

\subsection{相关性分析与滞后相关性分析}

首先采用 Pearson 相关系数和 Spearman 相关系数分析各变量与出厂水浊度之间的线性和单调关系。对于变量 $x_i$ 与目标变量 $y$，Pearson 相关系数为
\begin{equation}
  r_{x_i,y}=\frac{\sum_{t=1}^{n}(x_{i,t}-\bar{x}_i)(y_t-\bar{y})}
  {\sqrt{\sum_{t=1}^{n}(x_{i,t}-\bar{x}_i)^2}\sqrt{\sum_{t=1}^{n}(y_t-\bar{y})^2}} .
\end{equation}

进一步地，对不同滞后阶数下的变量与目标变量计算相关性，以识别潜在时滞影响。

\subsection{基于 XGBoost 的出厂水浊度预测模型}

为刻画输入变量与出厂水浊度之间的非线性关系，本文建立 XGBoost 回归模型。模型输入为当前变量及其滞后特征，输出为出厂水浊度 $y_t$。模型预测形式为
\begin{equation}
  \hat{y}_t=F\left(X_t,X_{t-1},\ldots,X_{t-12}\right).
\end{equation}

\subsection{基于 SHAP 的变量重要性解释}

为避免机器学习模型成为黑箱，本文采用 SHAP 方法解释不同变量对模型输出的贡献。SHAP 值为正表示该变量在对应样本中推高预测浊度，SHAP 值为负表示该变量降低预测浊度。

\subsection{模型性能与指定日期预测结果}

\begin{table}[H]
\centering
\caption{不同模型对出厂水 NTU 的预测性能对比}
\begin{tabular}{lccc}
\toprule
模型 & $RMSE$ & $MAE$ & $R^2$ \\
\midrule
多元线性回归 &  &  &  \\
随机森林 &  &  &  \\
XGBoost &  &  &  \\
GRU &  &  &  \\
\bottomrule
\end{tabular}
\end{table}

\begin{table}[H]
\centering
\caption{2026 年 2 月指定日期出厂水 NTU 预测结果}
\begin{tabular}{ccccc}
\toprule
日期 & 时间 & 预测值 & 真实值 & 绝对误差 \\
\midrule
2026-02-01 &  &  &  &  \\
2026-02-10 &  &  &  &  \\
2026-02-20 &  &  &  &  \\
\bottomrule
\end{tabular}
\end{table}

\section{问题二：滤后水浊度的时滞动态模型}

\subsection{滤后水浊度影响机制分析}

滤后水浊度 $FILT.NTU$ 是衡量混凝、沉淀和过滤效果的重要指标。原水浊度反映进入处理系统的颗粒负荷，原水 pH 会影响混凝剂效果，矾投加量是调控混凝反应的关键变量，原水流量则影响反应时间和过滤负荷。上述变量对滤后水浊度的影响并非即时发生，而是存在不同程度的时间滞后。

\subsection{动态时滞模型建立}

设 $z_t$ 表示 $t$ 时刻滤后水浊度，$x_{i,t-\tau_i}$ 表示第 $i$ 个输入变量在滞后 $\tau_i$ 后对滤后水浊度产生影响，则动态模型可表示为
\begin{equation}
  z_t=f\left(x_{1,t-\tau_1},x_{2,t-\tau_2},\ldots,x_{m,t-\tau_m}\right)+\varepsilon_t .
\end{equation}

其中，$\tau_i\in\{0,1,\ldots,12\}$，对应 0--24 小时滞后。

\subsection{最优时滞参数识别}

本文结合滞后相关性、网格搜索和模型特征重要性确定各变量的最优时滞。对于每个变量，在候选滞后阶数集合中选择使模型预测误差最小或特征贡献最大的滞后阶数。

\begin{table}[H]
\centering
\caption{主要输入变量的最优时滞参数}
\begin{tabular}{lccc}
\toprule
输入变量 & 最优滞后阶数 & 滞后时间/h & 作用方向 \\
\midrule
$R/W\ NTU$ &  &  &  \\
$R/W\ PH$ &  &  &  \\
$ALUM$ &  &  &  \\
$F/RIDE$ &  &  &  \\
$R/W\ FLOW$ &  &  &  \\
\bottomrule
\end{tabular}
\end{table}

\subsection{模型验证}

\begin{table}[H]
\centering
\caption{$FILT.NTU$ 动态模型预测性能}
\begin{tabular}{lccc}
\toprule
模型 & $RMSE$ & $MAE$ & $R^2$ \\
\midrule
ARX 基准模型 &  &  &  \\
滞后随机森林模型 &  &  &  \\
Lag-XGBoost 模型 &  &  &  \\
\bottomrule
\end{tabular}
\end{table}

\section{问题三：基于水力停留时间与 GRU 的出厂水浊度多步预测模型}
由于 18ML LEVEL 和 18ML FLOW 在本数据集中完全缺失，本文不将其纳入模型，而采用 C/W WELL LEVEL 与 T/W FLOW 构造近似水力停留时间特征。清水池具有调蓄和混合作用，当前出厂水浊度通常受到过去一段时间滤后水和清水池状态的共同影响。若用 \(L_t\) 表示清水池水位，用 \(Q_t\) 表示出厂水流量，则近似水力停留时间特征定义为
\begin{equation}
  HRT_t=\frac{L_t}{Q_t+\epsilon},
\end{equation}
其中 \(\epsilon\) 为防止分母为零的极小常数。
\subsection{水力停留时间特征构造}

清水池具有调蓄和混合作用，当前出厂水浊度通常受到过去一段时间滤后水和清水池状态的共同影响。若用 $L_t$ 表示清水池水位，用 $Q_t$ 表示出厂水流量，则近似水力停留时间特征定义为
\begin{equation}
  HRT_t=\frac{L_t}{Q_t+\epsilon},
\end{equation}
其中 $\epsilon$ 为防止分母为零的极小常数。

\subsection{多步预测问题定义}

由于采样间隔为 2 小时，未来 1--12 小时预测可转化为未来 1--6 步预测。设 $h=1,2,\ldots,6$，则多步预测目标为
\begin{equation}
  \hat{y}_{t+h}=g\left(y_{t-k:t},z_{t-k:t},X_{t-k:t},HRT_{t-k:t}\right),
\end{equation}
其中 $k$ 表示历史窗口长度。

\subsection{GRU 模型结构}

GRU 通过更新门和重置门刻画历史序列信息，其基本结构为
\begin{align}
  r_t &= \sigma(W_r x_t+U_r h_{t-1}+b_r),\\
  u_t &= \sigma(W_u x_t+U_u h_{t-1}+b_u),\\
  \tilde{h}_t &= \tanh(W_h x_t+U_h(r_t\odot h_{t-1})+b_h),\\
  h_t &= (1-u_t)\odot h_{t-1}+u_t\odot \tilde{h}_t .
\end{align}

\subsection{多步预测性能评价}

\begin{table}[H]
\centering
\caption{不同预测步长下模型预测性能}
\begin{tabular}{ccccc}
\toprule
预测步长 & 对应时间 & $RMSE$ & $MAE$ & $R^2$ \\
\midrule
1 step & 2 h &  &  &  \\
2 steps & 4 h &  &  &  \\
3 steps & 6 h &  &  &  \\
4 steps & 8 h &  &  &  \\
5 steps & 10 h &  &  &  \\
6 steps & 12 h &  &  &  \\
\bottomrule
\end{tabular}
\end{table}

\subsection{指定日期 7:00 至 19:00 预测结果}

\begin{table}[H]
\centering
\caption{2026 年 2 月指定日期 7:00 至 19:00 出厂水 NTU 预测结果}
\begin{tabular}{cccccc}
\toprule
日期 & 时间 & 2h预测 & 4h预测 & 6h预测 & 12h预测 \\
\midrule
2026-02-01 & 07:00 &  &  &  &  \\
2026-02-01 & 09:00 &  &  &  &  \\
2026-02-01 & 11:00 &  &  &  &  \\
2026-02-01 & 13:00 &  &  &  &  \\
2026-02-01 & 15:00 &  &  &  &  \\
2026-02-01 & 17:00 &  &  &  &  \\
2026-02-01 & 19:00 &  &  &  &  \\
\bottomrule
\end{tabular}
\end{table}

\subsection{输入变量敏感性分析}

为分析原水水质突变和矾投加调整对未来出厂水浊度预测结果的影响，本文设计情景扰动实验。设原始输入为 $X_t$，扰动后输入为 $X'_t$，则相对敏感度定义为
\begin{equation}
  S_i=\frac{\Delta \hat{y}/\hat{y}}{\Delta x_i/x_i} .
\end{equation}

\section{问题四：水质风险评价模型}

\subsection{风险评价指标构建}

以出厂水浊度 $y_t$ 为核心指标，根据生活饮用水浊度限值 $NTU\leq 1$ 定义超标幅度
\begin{equation}
  E_t=\max(0,y_t-1).
\end{equation}

同时定义浊度波动指标
\begin{equation}
  V_t=|y_t-y_{t-1}|.
\end{equation}

设 $D_t$ 表示连续异常持续时间，则综合风险指数为
\begin{equation}
  R_t=\alpha E_t+\beta D_t+\gamma V_t,
\end{equation}
其中 $\alpha,\beta,\gamma$ 为权重，满足 $\alpha+\beta+\gamma=1$。

\subsection{风险等级划分}

\begin{table}[H]
\centering
\caption{水质风险等级划分标准}
\begin{tabularx}{0.95\textwidth}{l L L}
\toprule
风险等级 & 判定依据 & 含义 \\
\midrule
安全 & 未超标且波动较小 & 水质稳定达标 \\
低风险 & 接近阈值或短时轻微波动 & 需要关注 \\
中风险 & 短时超标或波动明显 & 需要采取调控措施 \\
高风险 & 持续超标或严重异常 & 需要立即干预 \\
\bottomrule
\end{tabularx}
\end{table}

\subsection{2026 年近三个月风险评价结果}

\begin{table}[H]
\centering
\caption{2026 年 1--3 月各风险等级天数占比}
\begin{tabular}{lcc}
\toprule
风险等级 & 天数 & 占比 \\
\midrule
安全 &  &  \\
低风险 &  &  \\
中风险 &  &  \\
高风险 &  &  \\
\bottomrule
\end{tabular}
\end{table}

\begin{table}[H]
\centering
\caption{2026 年 3 月每日风险分类结果}
\begin{tabular}{ccccc}
\toprule
日期 & 日均 NTU & 日最大 NTU & 连续超标时长/h & 风险等级 \\
\midrule
2026-03-01 &  &  &  &  \\
2026-03-02 &  &  &  &  \\
\multicolumn{5}{c}{\kaishu 此处补充 3 月完整日期结果} \\
\bottomrule
\end{tabular}
\end{table}

\section{模型评价与敏感性分析}

本节从预测精度、解释性、稳定性和实际应用价值四个方面评价模型。预测精度采用 $RMSE$、$MAE$ 和 $R^2$ 衡量；解释性通过 SHAP 重要性、时滞参数和风险指标构成体现；稳定性通过不同月份、不同预测步长和输入扰动情景下的模型表现进行检验。

\section{模型优缺点}

\subsection{模型优点}
\begin{enumerate}[label=（\arabic*）]
  \item 充分考虑水处理过程的时间滞后特征，避免静态建模的局限；
  \item 将 XGBoost、SHAP 与 GRU 结合，兼顾预测精度和可解释性；
  \item 引入近似水力停留时间特征，使模型具有一定机理解释；
  \item 风险评价同时考虑超标幅度、持续时间和波动程度，较单一阈值判断更合理。
\end{enumerate}

\subsection{模型不足}
\begin{enumerate}[label=（\arabic*）]
  \item 清水池水力停留时间只能由水位和流量近似表征，可能与真实停留时间存在偏差；
  \item 部分变量缺失率较高，可能影响完整变量模型的稳定性；
  \item 对极端异常工况的预测能力依赖历史样本覆盖情况。
\end{enumerate}

\section{结论}

本文围绕自来水厂水浊度预测与风险评估问题，构建了从数据预处理、影响因素筛选、时滞动态建模、多步预测到风险分级的完整建模框架。后续将在完成数据清洗和模型训练后，补充各问题的具体数值结果、预测表格和图形分析。

\newpage

% =========================================================
% 参考文献
% =========================================================

\begin{thebibliography}{99}

\bibitem{gb5749}
国家市场监督管理总局，中国国家标准化管理委员会．生活饮用水卫生标准：GB 5749--2022[S]．北京：中国标准出版社，2022．

\bibitem{qasim}
Qasim S R, Motley E M, Zhu G. Water Works Engineering: Planning, Design and Operation[M]. Upper Saddle River: Prentice Hall, 2000.

\bibitem{box}
Box G E P, Jenkins G M, Reinsel G C, Ljung G M. Time Series Analysis: Forecasting and Control[M]. 5th ed. Hoboken: Wiley, 2015.

\bibitem{breiman}
Breiman L. Random Forests[J]. Machine Learning, 2001, 45(1): 5--32.

\bibitem{xgboost}
Chen T, Guestrin C. XGBoost: A Scalable Tree Boosting System[C]. Proceedings of the 22nd ACM SIGKDD International Conference on Knowledge Discovery and Data Mining, 2016: 785--794.

\bibitem{shap}
Lundberg S M, Lee S I. A Unified Approach to Interpreting Model Predictions[C]. Advances in Neural Information Processing Systems, 2017, 30: 4765--4774.

\bibitem{gru}
Cho K, van Merrienboer B, Gulcehre C, Bahdanau D, Bougares F, Schwenk H, Bengio Y. Learning Phrase Representations using RNN Encoder--Decoder for Statistical Machine Translation[C]. Proceedings of EMNLP, 2014: 1724--1734.

\bibitem{lstm}
Hochreiter S, Schmidhuber J. Long Short-Term Memory[J]. Neural Computation, 1997, 9(8): 1735--1780.

\bibitem{hyndman}
Hyndman R J, Athanasopoulos G. Forecasting: Principles and Practice[M/OL]. 3rd ed. Melbourne: OTexts, 2021. \url{https://otexts.com/fpp3/}，访问时间：2026年6月12日.

\bibitem{goodfellow}
Goodfellow I, Bengio Y, Courville A. Deep Learning[M]. Cambridge: MIT Press, 2016.

\end{thebibliography}

\newpage

\end{document}